\documentclass[a4paper]{article}

\usepackage{a4wide}
\usepackage[utf8]{inputenc}
\usepackage[ngerman]{babel}
\usepackage{fancyhdr}
\usepackage{graphicx}
\usepackage[parfill]{parskip}
\usepackage{color}

% fuer "deko-elemente"
\pagestyle{fancy}

% we needs moar platz
\topmargin -5mm
\headheight 20mm

% include logo in document head
\chead{\includegraphics[width=50mm]{logo.pdf}}

% seitenzahl entfernen
\cfoot{}

% footer content
\lfoot{\small{AlphaLabs}}
\rfoot{\tiny{Alexander Drangmeister, Dennis Esders, Philipp Heisig, Robert Baruck, Stephan Wilde, Kilian Költzsch}}

% headrule
\renewcommand\headrule{
\vspace{3mm}
\footrule
}

% footrule
\renewcommand\footnoterule{\rule{\linewidth}{0pt}}

% change standard font
\usepackage[T1]{fontenc}
\newcommand{\changefont}[3]{
\fontfamily{#1} \fontseries{#2} \fontshape{#3} \selectfont}

% Tagesordnungspunkt Item
\newcommand{\TOP}[1]{\item \textbf{#1}\par}



\begin{document}
\changefont{cmss}{m}{n} %computer modern sans serif

\begin{center}
\textbf{\Large Protokoll zum Gruppentreffen vom 27. April 2014}
\end{center}
\vspace{3mm}

\textbf{Beginn:} 13:00 Uhr, \textbf{Ende:} 14:00 Uhr\\
\textbf{Anwesende:} Alex, Dennis, Kilian

\vspace{5mm}

\begin{enumerate}

\TOP{Besprechung der nächsten Abgabe}
Wir gehen die benötigten Teilbereiche für die nächste Abgabe durch und besprechen, was kurzfristig noch erledigt werden muss.

\textbf{Ziel}\\
Genauere Details sowohl zur Zielgruppe und Zielstellung sind bereits ausgearbeitet.

\textbf{Corporate Design}\\
Das Design der Dokumente steht soweit und ist anhand dieses Protokolls ersichtlich. Für umfangreichere Dokumente sind jedoch Änderungen an Details möglich.

\textbf{Arbeitspakete}\\
Der Ablauf des Projekts wurde in drei große Teilbereiche gegliedert.

\textit{Design}
\begin{itemize}
\item Spielelemente
\item Interface
\end{itemize}

\textit{Entwicklung}
\begin{itemize}
\item Engine/Framework
\item Programmierung
\end{itemize}

\textit{Inhalte}
\begin{itemize}
\item Story
\item Recherche
\item Minispielinhalte
\end{itemize}


\textbf{Zeitplan}\\
Ein Gantt Diagramm für einen übersichtlichen Zeitplan ist noch am Entstehen und wird im Dokument zur Teilzielabgabe enthalten sein.

\textbf{Inhalte-Aufbau}\\
Genauere Details zu Inhalten und bereits bestehende Mockups werden ebenfalls zur nächsten Teilzielabgabe mit festgehalten.

\end{enumerate}


% das ist herrlich kompliziert, aber es klappt! das ist alles was zaehlt! :D
\definecolor{dgreen}{rgb}{0.549,0.776,0.247}
\makeatletter
\def\footrule{{
  \vskip-\footruleskip\vskip-\footrulewidth
  \color{\footrulecolor}
  \hrule\@width\headwidth\@height
  \footrulewidth\vskip\footruleskip
}}
\makeatother
\renewcommand{\footrulewidth}{3pt}
\newcommand{\footrulecolor}{dgreen}

\end{document}
